
\def\version{1}
\problemname{Prominence}

``What is the world's tallest mountain? The highest point on Mount Everest's summit. OK, but what is the world's second tallest mountain? The second highest point on Mount Everest's summit, of course.''\\
With this logic, the list of the world's tallest mountains becomes very silly. But there is a solution: one introduces the concept of prominence. The prominence of a mountain is the minimum vertical distance one must descend from the mountain to reach a strictly higher mountain. It works as a kind of measure of how independent a mountain is, and if one removes all points with prominence less than 200 m, one gets rid of all silly small mountains that are really sitting on higher mountains. This problem is about finding all prominences in a graph.
\\\\
We have a graph with $n$ nodes and $m$ edges, where each node $i$ has a non-negative integer $h(i)$, the node's height. A node's prominence $P(i)$ is the minimum height one must descend from the node to reach a node with strictly higher height. Here is a slightly more mathematical definition:
Let $G(i)$ be the set of all paths from node $i$ to some other node $j$ such that $h(j) > h(i)$. The prominence of $i$ is defined as

\begin{align*} 
P(i) =  \min_{\gamma \in G(i)}  \left\{   h(i) - \min_{k\in \gamma}(h(k))  \right\}
\end{align*} 

If $G(i) = \emptyset$, i.e.\ if it is not at all possible to reach a node with higher height, then we say that the prominence is $h(i)$. \\\\
Given a graph, find all nodes' prominences.

\section*{Input}
One line with two integers, $n$ and $m$. One line with $n$ integers $0 \leq h(i) \leq 10^9$, the nodes' heights. $m$ lines with two integers, $a$ and $b$ ($1 \leq a,b \leq n$), meaning that an edge goes between nodes $a$ and $b$.

\section*{Output}
One line with $n$ integers, the nodes' prominences.

\section*{Scoring}
Your solution will be tested on a set of test groups. To earn points for a group, you must pass all test cases in that group.

\noindent
\begin{tabular}{| l | l | l | l |}
\hline
Group & Points & Constraints \\ \hline
	1     & 25 & $n \le 1000, m \le 4000$ \\ \hline
	2     & 25 & $n \le 100\,000,$ the graph is a line, i.e.\ the edges are $(1,2),(2,3) \dots (n-1,n)$  \\ \hline
	3     & 50 & $n \le 100\,000, m \le 400\,000$ \\ \hline

\end{tabular}
